\section{Statistical methodology and full contrasts}
\label{app:stats}

\paragraph{What the intervals describe.}
Every Protocol~A result is evaluated three times, with decoding seeds 42, 1234, and 2025,
and reported as a mean with a sample standard deviation over those three runs. Under
greedy decoding the residual run-to-run variation comes from batching and kernel
nondeterminism in the inference engine rather than from sampling, which is why the
standard deviations are small but nonzero, and occasionally large on the 30-problem
contest-mathematics sets where a single problem is worth 3.3 points. These seeds do not
estimate variability across independent pre-training runs. Doing so would require
repeating the 300B-token pre-training runs, which was out of reach.

\paragraph{Paired item-level bootstrap.}
Contrasts on the suite average carry 95\% percentile intervals from a paired item-level
bootstrap with 2{,}000 replicates. Scores are averaged over the three seeds first, giving
one score per item per checkpoint. Each replicate resamples items with replacement
\emph{within} each benchmark, which preserves the composition of the eleven-benchmark
macro average in every replicate, and the contrast is recomputed on the resampled items.
Because the two checkpoints of a contrast are evaluated on the same items, the resampling is
paired and the interval reflects the uncertainty of the difference rather than of the two
levels separately. JEEBench awards partial credit, so item scores lie in $[0,1]$ and are
treated as such rather than as Bernoulli outcomes. Two-sided $p$-values are taken from
the bootstrap distribution and corrected within each contrast family with the
Benjamini-Hochberg procedure. The three families are the front-loading effects, the
substrate contrasts at matched front-loading, and the interaction, and
Table~\ref{tab:contrasts} reports all of them.

% Generated by generate_workshop_artifacts.py. Do not edit by hand.
\begin{table}[t]
\centering
\caption{All Protocol~A contrasts on the eleven-benchmark suite average, with 95\% paired item-level bootstrap intervals (2{,}000 replicates, items resampled within benchmark, decoding seeds averaged first) and Benjamini--Hochberg adjusted $q$-values computed within each family. Positive values favour front-loading in the first block and the permissive substrate in the second. Estimates are computed on unrounded item-level scores, so they can differ in the final digit from the difference of the rounded values printed in Table~\ref{tab:factorial}.}
\label{tab:contrasts}
\small
\setlength{\tabcolsep}{5pt}
\begin{tabular}{llrlr}
\toprule
\textbf{Contrast} & \textbf{Stage} & \textbf{Estimate} & \textbf{95\% CI} & \textbf{$q$} \\
\midrule
\multicolumn{5}{l}{\emph{Front-loading effect, within web substrate, 4K context}} \\
\quad FineWeb-Edu & Base & $+19.12$ & $[+17.25, +20.88]$ & $<0.001$ \\
\quad FineWeb-Edu & SFT & $+15.57$ & $[+14.31, +17.06]$ & $<0.001$ \\
\quad FineWeb-Edu & DPO & $+15.87$ & $[+14.42, +17.36]$ & $<0.001$ \\
\quad MixtureVitae & Base & $+10.27$ & $[+8.90, +11.67]$ & $<0.001$ \\
\quad MixtureVitae & SFT & $+11.82$ & $[+10.27, +13.41]$ & $<0.001$ \\
\quad MixtureVitae & DPO & $+12.44$ & $[+10.73, +14.20]$ & $<0.001$ \\
\quad Nemotron-CC-HQ & Base & $+19.17$ & $[+17.38, +21.00]$ & $<0.001$ \\
\quad Nemotron-CC-HQ & SFT & $+14.45$ & $[+13.05, +15.89]$ & $<0.001$ \\
\quad Nemotron-CC-HQ & DPO & $+15.23$ & $[+13.81, +16.63]$ & $<0.001$ \\
\addlinespace[0.2em]
\multicolumn{5}{l}{\emph{Front-loading effect, within web substrate, 16K context}} \\
\quad FineWeb-Edu & Base & $+19.06$ & $[+17.43, +20.77]$ & $<0.001$ \\
\quad FineWeb-Edu & SFT & $+14.70$ & $[+13.36, +16.08]$ & $<0.001$ \\
\quad FineWeb-Edu & DPO & $+13.63$ & $[+12.23, +15.05]$ & $<0.001$ \\
\quad MixtureVitae & Base & $+19.09$ & $[+17.34, +20.82]$ & $<0.001$ \\
\quad MixtureVitae & SFT & $+15.55$ & $[+13.71, +17.42]$ & $<0.001$ \\
\quad MixtureVitae & DPO & $+15.76$ & $[+13.94, +17.62]$ & $<0.001$ \\
\quad Nemotron-CC-HQ & Base & $+19.11$ & $[+17.49, +20.83]$ & $<0.001$ \\
\quad Nemotron-CC-HQ & SFT & $+13.95$ & $[+12.68, +15.29]$ & $<0.001$ \\
\quad Nemotron-CC-HQ & DPO & $+14.07$ & $[+12.72, +15.48]$ & $<0.001$ \\
\midrule
\multicolumn{5}{l}{\emph{Permissive minus non-permissive substrate, both front-loaded, 4K context}} \\
\quad MV - FineWeb-Edu & Base & $-7.49$ & $[-9.30, -5.59]$ & $<0.001$ \\
\quad MV - FineWeb-Edu & SFT & $+0.34$ & $[-1.28, +1.91]$ & $0.656$ \\
\quad MV - FineWeb-Edu & DPO & $+0.94$ & $[-0.90, +2.80]$ & $0.429$ \\
\quad MV - Nemotron & Base & $-5.39$ & $[-7.36, -3.41]$ & $<0.001$ \\
\quad MV - Nemotron & SFT & $+0.50$ & $[-1.18, +2.26]$ & $0.645$ \\
\quad MV - Nemotron & DPO & $+0.40$ & $[-1.20, +2.01]$ & $0.656$ \\
\addlinespace[0.2em]
\multicolumn{5}{l}{\emph{Permissive minus non-permissive substrate, both front-loaded, 16K context}} \\
\quad MV - FineWeb-Edu & Base & $+3.54$ & $[+1.75, +5.39]$ & $0.002$ \\
\quad MV - FineWeb-Edu & SFT & $+4.77$ & $[+3.04, +6.60]$ & $<0.001$ \\
\quad MV - FineWeb-Edu & DPO & $+5.66$ & $[+3.94, +7.47]$ & $<0.001$ \\
\quad MV - Nemotron & Base & $+3.41$ & $[+1.63, +5.10]$ & $<0.001$ \\
\quad MV - Nemotron & SFT & $+5.05$ & $[+3.21, +6.96]$ & $<0.001$ \\
\quad MV - Nemotron & DPO & $+5.08$ & $[+3.12, +7.01]$ & $<0.001$ \\
\midrule
\multicolumn{5}{l}{\emph{Interaction: front-loading effect, permissive minus non-permissive, 4K context}} \\
\quad MV - non-permissive & Base & $-8.87$ & $[-10.67, -7.01]$ & $0.002$ \\
\quad MV - non-permissive & SFT & $-3.19$ & $[-4.97, -1.46]$ & $0.002$ \\
\quad MV - non-permissive & DPO & $-3.11$ & $[-4.94, -1.20]$ & $0.003$ \\
\addlinespace[0.2em]
\multicolumn{5}{l}{\emph{Interaction: front-loading effect, permissive minus non-permissive, 16K context}} \\
\quad MV - non-permissive & Base & $+0.00$ & $[-1.99, +1.95]$ & $0.974$ \\
\quad MV - non-permissive & SFT & $+1.22$ & $[-0.81, +3.23]$ & $0.275$ \\
\quad MV - non-permissive & DPO & $+1.91$ & $[-0.11, +3.98]$ & $0.103$ \\
\bottomrule
\end{tabular}
\end{table}


\paragraph{Robustness to how the suite is weighted.}
The eleven-task macro average contains five mathematics and three code benchmarks, so it
implicitly weights those domains at 45\% and 27\%. Table~\ref{tab:domain_robustness}
recomputes the two headline quantities under an equal-domain average, which gives each of
the four capability domains equal weight, and with each domain removed in turn. Neither
the front-loading effect nor the substrate contrast depends on the weighting.

% Generated by generate_workshop_artifacts.py. Do not edit by hand.
\begin{table}[t]
\centering
\caption{Robustness of the two headline quantities to how the suite average is weighted, at 16K context after Tulu3 SFT. The eleven-task macro average implicitly weights math at 45\% and code at 27\%, and neither the front-loading effect nor the substrate contrast depends on that weighting. Protocol~A (Appendix Table~\ref{tab:eval-settings-A}). Values are mean$_{\pm\text{SD}}$ over decoding seeds 42, 1234, and 2025.}
\label{tab:domain_robustness}
\small
\begin{tabular}{lrrrrr}
\toprule
& \multicolumn{3}{c}{\textbf{Front-loading effect}} & \multicolumn{2}{c}{\textbf{Permissive $-$ non-permissive}} \\
\cmidrule(lr){2-4}\cmidrule(lr){5-6}
\textbf{Averaging view} & MV & Nemotron & FineWeb & vs.\ Nemotron & vs.\ FineWeb \\
\midrule
Suite average (11 tasks) & $+15.5$ & $+13.9$ & $+14.7$ & $+5.0$ & $+4.8$ \\
Equal-domain average & $+12.2$ & $+12.8$ & $+13.6$ & $+3.2$ & $+3.4$ \\
Drop math domain & $+11.7$ & $+15.0$ & $+15.4$ & $+3.7$ & $+4.2$ \\
Drop code domain & $+13.6$ & $+10.3$ & $+11.1$ & $+3.6$ & $+3.0$ \\
Drop science domain & $+19.0$ & $+16.1$ & $+17.3$ & $+6.5$ & $+6.2$ \\
Drop instruction domain & $+16.3$ & $+14.3$ & $+14.9$ & $+5.7$ & $+5.2$ \\
\bottomrule
\end{tabular}
\end{table}


\paragraph{The contest-mathematics benchmarks.}
AIME24 and AIME25 contain 30 problems each and AMC23 contains 40. At those sizes a paired
binomial power analysis puts the effect needed for 80\% power at roughly 28 points for the
AIME sets and 30 points for AMC23, so their percentage-point estimates are intrinsically
imprecise and we do not build effect-size claims on them. Where we quote them, as in the
pre-training ablation of Section~\ref{sec:receptivity} where AIME24 and AIME25 fall to
exactly zero, the claim is that a model does or does not clear the benchmark floor, which
is a qualitatively different and far more robust statement than an estimate of how far it
clears it by.
